% Course: 77972 -- Physical Models of Biological Systems
% Academic year: 2026--2027 | Semester: A
% Actual lecture: 3 | Document chapter: 3
% Lecture date: 2026-11-03
% Title: הסקה מנתונים, התפלגויות בדידות וניסוי לוריא--דלברוק
% Lecturer: Dr. Amir Erez
% Summarizer: Unknown (AI-assisted preparation)
% Prepared: 2026-11-05
% Editorial revision: 2026-11-05 -- Hebrew prose and explanatory transitions.
% Lecture_Summaries/Lec03_Discrete_Distributions_and_Luria_Delbruck.tex (secondary, checked critically).
% partition noise through 01:35; Poisson and fluctuation test thereafter;
% PGF 02:14 onward; moments explicitly deferred at 02:37.
% Editorial choice: generations n=1,...,g consistently; no initial resistant cells.
% Supplementary examples and completion of variance are identified in the text.

\chapter{הסקה מנתונים, התפלגויות בדידות וניסוי לוריא--דלברוק}\label{chap:he03}
\noindent\textbf{הרצאה 3 \enspace|\enspace 3 בנובמבר 2026}

במדידות ביולוגיות נתקלים בהבדלים בין תאים, ברעש בחלוקת מולקולות ובמוטציות נדירות שצאצאיהן מתרבים. בהרצאה זו משתמשים בהסתברות כדי להסיק השערות מנתונים. בהמשך נראה כיצד תנודות בין תאים מכיילות מדידת פלואורסצנציה, וכיצד תנודות בין תרביות מלמדות מתי נוצרת עמידות בחיידקים.

\subsection*{מטרות הלמידה}
\begin{itemize}
\item להבחין בין נראות, הסתברות מוקדמת והסתברות מעודכנת, ולפתור את דוגמת הבדיקה הרפואית.
\item לפרש מתאם ושונות משותפת ולהסביר מדוע יש לצייר את הנתונים עצמם.
\item לגזור את התוחלת והשונות של ספירות בינומיות ופואסוניות ולכייל פלואורסצנציה מרעש החלוקה.
\item לבנות מודל למוטציות בזמן גידול, לחשב הסתברות לאפס עמידים, ולאמוד ממנה את פרמטר המוטציה.
\end{itemize}

\section{משפט בייס: מהנתונים אל ההשערה}\label{sec:he03-bayes}
בביולוגיה יש לעיתים כמה מנגנונים אפשריים לאותה תופעה. לכן השאלה אינה רק מה מודל נתון מנבא, אלא עד כמה התצפית תומכת בו ביחס לחלופות. נסמן השערה ב-$H$ ותצפית ב-$E$.

\begin{definition}{מרכיבי ההסקה הבייזיאנית}{he03-bayes}
ההסתברות המוקדמת $P(H)$ מתארת את הידע לפני התצפית. הנראות $P(E\mid H)$ מתארת את הסתברות התצפית בהינתן ההשערה; לאחר קיבוע הנתונים היא פונקציה של ההשערה. ההסתברות המעודכנת $P(H\mid E)$ מתארת את הידע לאחר התצפית. המכנה $P(E)$ הוא הסתברות התצפית הכוללת, המשמשת לנרמול.
\end{definition}

\begin{theorem}{משפט בייס}{he03-bayes-rule}
אם $P(E)>0$, אז
\begin{equation}\label{eq:he03-bayes}
P(H\mid E)=\frac{P(E\mid H)P(H)}{P(E)}.
\end{equation}
עבור השערות זרות וממצות $H_j$, המכנה הוא
\begin{equation}\label{eq:he03-evidence}
P(E)=\sum_j P(E\mid H_j)P(H_j).
\end{equation}
\end{theorem}
\begin{proof}
כותבים את ההסתברות המשותפת בשתי דרכים:
\begin{equation}\label{eq:he03-joint}
P(H,E)=P(E\mid H)P(H)=P(H\mid E)P(E).
\end{equation}
חלוקה ב-$P(E)$ נותנת את התוצאה. סכימה על חלוקה ממצה של מרחב ההשערות נותנת את המכנה.
\end{proof}

נראות אינה התפלגות הסתברות על ההשערות. בפרט, מקסום נראות אינו בהכרח מקסום ההסתברות המעודכנת: האחרון כולל גם את ההסתברות המוקדמת. בחירה אחידה אינה מבטלת את בחירת המודל. למשל, צפיפות אחידה בפרמטר חיובי $\theta$ אינה אחידה בלוגריתם שלו: עבור $u=\ln\theta$, כלל החלפת המשתנים נותן
\begin{equation}\label{eq:he03-prior-transform}
p_u(u)=p_\theta(\ee^u)\ee^u.
\end{equation}
החלפת המשתנה מבהירה מדוע בחירה אחידה מחייבת לציין באיזה פרמטר בוחרים אותה. הסתברות מוקדמת השווה לאפס משאירה את ההשערה בהסתברות אפס לאחר עדכון בעל מכנה חיובי; הסתברות קטנה אך חיובית יכולה לגדול בעקבות ראיות חזקות.

\subsection{בדיקה רפואית ושכיחות המחלה}\label{sec:he03-test}
נסמן מחלה ב-$D$, בריאות ב-$\bar D$, ותוצאה חיובית ב-$+$. השכיחות היא $\pi=P(D)$; הרגישות היא $a=P(+\mid D)$; הסגוליות היא $b=P(-\mid\bar D)$. מכאן שיעור החיוביים השגויים הוא $1-b$.
\begin{align}
P(+)&=a\pi+(1-b)(1-\pi),\label{eq:he03-positive}\\
P(D\mid+)&=\frac{a\pi}{a\pi+(1-b)(1-\pi)}.\label{eq:he03-ppv}
\end{align}

\begin{example}{מספרי דוגמת ההרצאה}{he03-test}
עבור $\pi=0.009$ ורגישות וסגוליות $a=b=0.97$ מתקבל
\begin{align}
P(+)&=0.97\cdot0.009+0.03\cdot0.991=0.03846,\label{eq:he03-test-total}\\
P(D\mid+)&=\frac{0.00873}{0.03846}\simeq0.227.\label{eq:he03-test-posterior}
\end{align}
כלומר, ההסתברות למחלה לאחר תוצאה חיובית היא כ-$23\%$, ולא $97\%$. כדי לראות מדוע, נחשוב על $100{,}000$ נבדקים: בתוחלת יהיו $900$ חולים, מהם $873$ חיוביים, ומבין $99{,}100$ הבריאים יהיו $2{,}973$ חיוביים שגויים. רוב החיוביים מגיעים אפוא מקבוצת הבריאים הגדולה. למרות זאת, הבדיקה הגדילה את ההסתברות למחלה בערך פי $25$ ביחס לשכיחות ההתחלתית. אלו מספרי הדוגמה הלימודית, ולא נתוני ביצועים רפואיים עדכניים.
\end{example}

\begin{remark}
רגישות וסגוליות הן הסתברויות מותנות שונות; תמלול שמחליף ביניהן אינו משנה את הגדרתן. עדכון חוזר אפשרי, אך הכפלת הנראויות של שתי בדיקות מחייבת אי-תלות שלהן \emph{בהינתן מצב המחלה}. מקור שגיאה משותף עלול לגרום לתוצאות תלויות, ולכן שתי תוצאות חיוביות אינן בהכרח שתי ראיות בלתי-תלויות.
\end{remark}

\section{שונות משותפת, מתאם והצגת הנתונים}\label{sec:he03-correlation}
כשמודדים שני חלבונים באותם תאים, רוצים לדעת אם סטיותיהם מהממוצע משתנות יחד. יש להגדיר על פני מה ממוצעים: תאים, בעלי חיים, זמנים או גנים. החלפת ציר הדגימה משנה את השאלה הביולוגית.

\begin{definition}{שונות משותפת ומתאם פירסון}{he03-cov}
למשתנים בעלי מומנט שני סופי, עם ממוצעים $\mu_X,\mu_Y$ וסטיות תקן חיוביות $\sigma_X,\sigma_Y$, נגדיר
\begin{align}
\operatorname{Cov}(X,Y)&=\mathbb{E}[(X-\mu_X)(Y-\mu_Y)]
=\mathbb{E}[XY]-\mu_X\mu_Y,\label{eq:he03-cov}\\
\rho_{XY}&=\frac{\operatorname{Cov}(X,Y)}{\sigma_X\sigma_Y}.\label{eq:he03-rho}
\end{align}
\end{definition}
לשונות המשותפת יחידות של מכפלת שני המשתנים; המתאם חסר יחידות. אי-שוויון קושי--שוורץ נותן $|\rho|\le1$. שוויון $|\rho|=1$ משמעו קשר אפיני מדויק בהסתברות אחת, כאשר השונויות אינן אפס. מתאם חיובי מתאר נטייה ליניארית לסטות יחד, ומתאם שלילי נטייה לסטות בכיוונים הפוכים; הוא אינו קובע סיבתיות.

\begin{proposition}{פעולות ליניאריות ושונות}{he03-linear}
לקבועים $a,b,c,d$,
\begin{align}
\operatorname{Cov}(aX+b,cY+d)&=ac\operatorname{Cov}(X,Y),\label{eq:he03-affine}\\
\operatorname{Var}(X+Y)&=\operatorname{Var}(X)+\operatorname{Var}(Y)
+2\operatorname{Cov}(X,Y).\label{eq:he03-sumvar}
\end{align}
\end{proposition}
\begin{proof}
מרכוז מסלק את $b,d$, והוצאת קבועים מהתוחלת נותנת את הנוסחה הראשונה. פיתוח ריבוע הסכום של המשתנים הממורכזים נותן את השנייה. לכן מותר לחבר שונויות של משתנים בלתי-תלויים, אך לא של כל שני משתנים.
\end{proof}

\begin{example}{אפס מתאם אינו אי-תלות: דוגמת השלמה}{he03-uncorrelated}
יהי $X$ אחיד ב-$[-1,1]$ ו-$Y=X^2$. מסימטריה $\mathbb{E}[X]=\mathbb{E}[X^3]=0$, ולכן
\begin{equation}\label{eq:he03-square-example}
\operatorname{Cov}(X,Y)=\mathbb{E}[X^3]-\mathbb{E}[X]\mathbb{E}[X^2]=0.
\end{equation}
ובכל זאת $Y$ נקבע לחלוטין על ידי $X$. אי-תלות גוררת אפס שונות משותפת כאשר המומנטים קיימים; הכיוון ההפוך אינו נכון באופן כללי.
\end{example}

\subsection{מתאם מדגמי, ערבול וספירמן}\label{sec:he03-spearman}
במדידות מזווגות $(x_i,y_i)$, עבור $i=1,\ldots,m$, אומדים את מתאם פירסון באמצעות
\begin{equation}\label{eq:he03-sample-pearson}
r=\frac{\sum_{i=1}^m(x_i-\bar x)(y_i-\bar y)}
{\sqrt{\sum_{i=1}^m(x_i-\bar x)^2\sum_{i=1}^m(y_i-\bar y)^2}}.
\end{equation}
הנוסחה דורשת פיזור לא אפסי בשני המשתנים. ערך כגון $r=0.2$ אינו מספיק כדי לקבוע אם הזיווג משמעותי: יש חשיבות למספר הדגימות ולמבנה הרעש. בהרצאה הוצע לערבל את תוויות $y_i$, להשאיר את ערכי כל משתנה כפי שהם, ולחשב מתאם לכל ערבול. כך הורסים את הזיווג המקורי ובונים התפלגות ייחוס תחת השערת האפס של היעדר קשר. משווים אליה את $r$ שנמדד. תקפות מבחן התמורות דורשת שהזיווגים אכן ניתנים להחלפה תחת השערת האפס; בדגימות תלויות בזמן או בקבוצות אין לערבל את כל הרשומות ללא התחשבות במבנה. תוצאה שאינה חריגה ביחס לערבולים אינה הוכחה לאי-תלות.

\begin{definition}{מתאם ספירמן}{he03-spearman}
מחליפים את ערכי כל משתנה בדרגותיהם במדגם, $R_i=\operatorname{rank}(x_i)$ ו-$S_i=\operatorname{rank}(y_i)$, ומחשבים מתאם פירסון בין הדרגות:
\begin{equation}\label{eq:he03-spearman}
r_S=\operatorname{Corr}(R,S).
\end{equation}
בערכים שווים מקצים בדרך כלל דרגה ממוצעת. נדרש פיזור לא אפסי של הדרגות.
\end{definition}
ספירמן מודד קשר מונוטוני גם כשהוא אינו ליניארי. ללא קשרי שוויון, אם סדר שתי הסדרות זהה מתקבל $r_S=1$, ואם הסדר הפוך מתקבל $r_S=-1$. המחיר הוא אובדן סקלת הערכים: הבדל זעיר והבדל עצום בין ערכים סמוכים אינם מובחנים בדרגות. לכן נקודה בעלת ערך גדול מאוד משפיעה פחות דרך גודלה המספרי, אך ספירמן עדיין אינו חסין מכל השפעת חריגים. קשר לא מונוטוני, כגון הפרבולה בדוגמה לעיל, יכול לחמוק גם ממנו. טענה בתמלול שלפיה היעדר קשר בדרגות פירושו היעדר קשר בכלל דורשת אפוא סייג זה. מומלץ להשוות פירסון וספירמן לצד ציור הנתונים, ולא לבחור מקדם רק מפני שערכו גדול יותר.

\subsection{מרחקים בענן נתונים ורביעיית אנסקומב}\label{sec:he03-anscombe}
לוקטור מדידות $\mathbf{x}$ נגדיר מטריצת שונות משותפת $\Sigma$. היא סימטרית וחיובית למחצה; איברי האלכסון הם השונויות. בקישור לדיון הגאומטרי בהרצאה, ענן מאורך יכול להיראות שונה לאחר סיבוב צירים אף שהמידע עצמו אינו משתנה. אם $\Sigma=U\Lambda U^T$ חיובית לחלוטין, אז בבסיס הצירים הראשיים $\mathbf{x}'=U^T(\mathbf{x}-\boldsymbol\mu)$ מרחק מהלנוביס הוא
\begin{equation}\label{eq:he03-mahalanobis}
d_M^2=(\mathbf{x}-\boldsymbol\mu)^T\Sigma^{-1}(\mathbf{x}-\boldsymbol\mu)
=\sum_i\frac{(x_i')^2}{\lambda_i}.
\end{equation}
המרחק מודד סטייה ביחידות הפיזור בכל כיוון, ולא מרחק גאומטרי גולמי. כיוון בעל שונות קטנה מקבל משקל גדול. הנוסחה מניחה מטריצה הפיכה; כיוון בעל שונות אפס דורש טיפול נפרד. זהו פירוט משלים של הקשר בין שונות, צירים ומרחקים.

רביעיית אנסקומב מציגה ארבע קבוצות נתונים בעלות ממוצעים, שונויות, מתאם וקו רגרסיה כמעט זהים, אך צורות שונות מאוד: ענן בקירוב ליניארי, קשר מעוקל, השפעה של נקודה חריגה, ומצב שבו נקודה יחידה קובעת כמעט לבדה את השיפוע. המסר הוא \textbf{לצייר את הנתונים לפני שמסיקים מסקנה מהמספר המסכם}. בפרט, מתאם לבדו אינו מספיק לזיהוי מנגנון או להצדקת מודל ליניארי.

\section{מניסויי ברנולי להתפלגות בינומית}\label{sec:he03-binomial}
במודל הפשוט ביותר כל מולקולה נשלחת לאחד משני תאי בת, וכל חיידק עשוי לעבור אירוע נדיר. נתחיל בניסיון יחיד.
\begin{definition}{ברנולי וספירה בינומית}{he03-binomial}
משתנה ברנולי $B$ מקבל $1$ בהסתברות $p$ ו-$0$ בהסתברות $1-p$. אם $B_1,\ldots,B_M$ בלתי-תלויים ובעלי אותו $p$, מספר ההצלחות $L=\sum_{i=1}^M B_i$ בינומי:
\begin{equation}\label{eq:he03-binomial}
P(L=\ell)=\binom{M}{\ell}p^\ell(1-p)^{M-\ell},\qquad \ell=0,\ldots,M.
\end{equation}
\end{definition}
יש $\binom{M}{\ell}$ דרכים לבחור את ההצלחות, וכל סדרה נתונה בעלת אותה הסתברות. משפט הבינום מבטיח נרמול:
\begin{equation}\label{eq:he03-bin-normalization}
\sum_{\ell=0}^M\binom{M}{\ell}p^\ell(1-p)^{M-\ell}=[p+(1-p)]^M=1.
\end{equation}

\begin{proposition}{תוחלת ושונות בינומית}{he03-bin-moments}
\begin{equation}\label{eq:he03-bin-moments}
\mathbb{E}[L]=Mp,\qquad \operatorname{Var}(L)=Mp(1-p).
\end{equation}
\end{proposition}
\begin{proof}
מכיוון ש-$B^2=B$, מתקיים $\mathbb{E}[B]=\mathbb{E}[B^2]=p$ ולכן $\operatorname{Var}(B)=p-p^2$. ליניאריות התוחלת ואי-תלות הניסיונות נותנות
\begin{align}
\mathbb{E}\!\left[\sum_iB_i\right]&=\sum_i p=Mp,\label{eq:he03-bin-mean-proof}\\
\operatorname{Var}\!\left(\sum_iB_i\right)&=\sum_i p(1-p)=Mp(1-p).\label{eq:he03-bin-var-proof}
\end{align}
\end{proof}
בדיקת קצה: ב-$p=0$ או $p=1$ אין אקראיות ולכן השונות אפס. ב-$p=1/2$ השונות היא $M/4$, וסטיית התקן היחסית בספירה היא $1/\sqrt M$. מספרים גדולים מקטינים רעש יחסי, אף שהרעש המוחלט גדל.

\section{הטריק של רוזנפלד: רעש החלוקה מכייל פלואורסצנציה}\label{sec:he03-rosenfeld}
עוצמת אור אינה מספר מולקולות. נניח שלאחר חיסור רקע, תא בעל $M$ מולקולות מסומנות נותן עוצמה $Y=\alpha M$. הפרמטר $\alpha$ הוא עוצמה למולקולה, והוא תלוי במערכת המדידה. כיצד נמדוד אותו בתוך הניסוי עצמו?

מתבוננים בסרט של חלוקת תא, ומשייכים את שתי הבנות לאותה אם. בהינתן $M_0$ מולקולות באם, נניח שכל מולקולה עוברת לכל בת בהסתברות $1/2$, באופן בלתי-תלוי במולקולות האחרות, ושאין ייצור, פירוק או אובדן ניכרים בחלון החלוקה. אז
\begin{equation}\label{eq:he03-partition}
M_1\mid M_0\sim\operatorname{Bin}(M_0,1/2),\qquad M_2=M_0-M_1.
\end{equation}
לשתי הבנות יש אנטי-מתאם בהינתן $M_0$: מולקולה שהגיעה לאחת אינה יכולה להגיע לשנייה. נגדיר את ההפרש החתום $\Delta M=M_1-M_2=2M_1-M_0$. לכן
\begin{align}
\mathbb{E}[\Delta M\mid M_0]&=2(M_0/2)-M_0=0,\label{eq:he03-delta-mean}\\
\operatorname{Var}(\Delta M\mid M_0)&=4\operatorname{Var}(M_1\mid M_0)
=4\frac{M_0}{4}=M_0.\label{eq:he03-delta-var}
\end{align}
השגיאה הנפוצה היא להתייחס ל-$M_1,M_2$ כבלתי-תלויים; הדבר היה נותן בטעות $M_0/2$ במקום $M_0$.

\begin{figure}[htbp]
\centering
\begin{english}
\begin{tikzpicture}[>=stealth, every node/.style={font=\small}]
\node[draw,rounded corners,fill=primary!8,minimum width=27mm,minimum height=10mm] (m) at (0,0) {$M_0,\;Y_0=\alpha M_0$};
\node[draw,rounded corners,fill=primary!8,minimum width=26mm,minimum height=10mm] (a) at (-2.4,-1.5) {$M_1,\;Y_1=\alpha M_1$};
\node[draw,rounded corners,fill=primary!8,minimum width=26mm,minimum height=10mm] (b) at (2.4,-1.5) {$M_2,\;Y_2=\alpha M_2$};
\draw[->] (m)--node[above left] {$p=1/2$}(a);
\draw[->] (m)--node[above right] {$p=1/2$}(b);
\node at (0,-2.5) {$M_1+M_2=M_0,\qquad \Delta Y=Y_1-Y_2$};
\end{tikzpicture}
\end{english}
\caption{סכמת חלוקה מקורית: מספר המולקולות הכולל נשמר, אך החלוקה בין הבנות רועשת. מודדים משפחות רבות של אם ושתי בנות, ומשווים משפחות בעלות עוצמת אם דומה.}\label{fig:he03-partition}
\end{figure}

בהנחת אותו מקדם $\alpha$ לאם ולבנות, $\Delta Y=\alpha\Delta M$ ו-$Y_0=\alpha M_0$. מכאן
\begin{equation}\label{eq:he03-calibration}
\boxed{\operatorname{Var}(\Delta Y\mid Y_0)=\alpha Y_0},\qquad
\boxed{\sigma_{\Delta Y\mid Y_0}=\sqrt{\alpha Y_0}}.
\end{equation}
מתאימים את השונות מול $Y_0$: השיפוע הוא $\alpha$. אם מציירים סטיית תקן מול $\sqrt{Y_0}$, השיפוע הוא $\sqrt\alpha$ ולכן יש לרבע אותו. מול $Y_0$ עצמו מקבלים עקומת שורש, ולא ישר. בכל תחום עוצמות נדרשים זוגות בת רבים כדי לאמוד שונות; זוג בודד אינו מספק סטיית תקן.

\begin{example}{כיול מספרי: המחשה משלימה}{he03-calibration}
אם עבור עוצמות אם סביב $Y_0=1000$ יחידות נמדדת שונות הפרש $4000$ יחידות עוצמה בריבוע, השיפוע המקומי במודל הוא $\alpha=4$ יחידות למולקולה, ולכן $M_0=250$. היחידות במשוואה נכונות: עוצמה למולקולה כפול עוצמה נותנות עוצמה בריבוע כאשר מספר המולקולות הוא ספירה חסרת ממד.
\end{example}

\begin{remark}
הכיול מותנה במודל: חלוקה לא סימטרית, צברים של חלבונים, רעש מצלמה, רקע משתנה, הלבנה או הבשלה של הסמן עלולים להוסיף שונות. למשל, רעשי מדידה עצמאיים בבנות מוסיפים את סכום שונויותיהם ל-$\operatorname{Var}(\Delta Y)$. התאמה לחוק החזוי תומכת במודל החלוקה, אך אינה מוכיחה באופן יחיד שכל הנחותיו נכונות.
\end{remark}

\section{גבול האירועים הנדירים: פואסון}\label{sec:he03-poisson}
נעבור מ-$M$ ניסיונות עם הסתברות $p$ לגבול $M\to\infty$, $p\to0$, כאשר התוחלת $Mp=\lambda$ קבועה. מספר ההצלחות עצמו אינו קבוע; הוא נשאר אקראי. עבור $\ell$ קבוע,
\begin{align}
P(L=\ell)&=\frac{M(M-1)\cdots(M-\ell+1)}{\ell!}
\left(\frac{\lambda}{M}\right)^\ell
\left(1-\frac{\lambda}{M}\right)^{M-\ell}\label{eq:he03-poisson-limit-start}\\
&=\frac{\lambda^\ell}{\ell!}\prod_{j=0}^{\ell-1}\left(1-\frac jM\right)
\left(1-\frac{\lambda}{M}\right)^M
\left(1-\frac{\lambda}{M}\right)^{-\ell}
\longrightarrow\frac{\ee^{-\lambda}\lambda^\ell}{\ell!}.\label{eq:he03-poisson-limit}
\end{align}
המכפלה והגורם האחרון שואפים ל-$1$, ואילו הגורם האמצעי שואף ל-$\ee^{-\lambda}$.

\begin{definition}{התפלגות פואסון ומקדם פאנו}{he03-poisson}
ל-$L\sim\operatorname{Pois}(\lambda)$, עם $\lambda\ge0$, פונקציית ההסתברות היא תוצאת הגבול לעיל, עבור $\ell=0,1,\ldots$. כאשר התוחלת חיובית, מקדם פאנו מוגדר
\begin{equation}\label{eq:he03-fano-definition}
F=\frac{\operatorname{Var}(L)}{\mathbb{E}[L]}.
\end{equation}
\end{definition}
הנרמול נובע מטור האקספוננט: $\ee^{-\lambda}\sum_{\ell\ge0}\lambda^\ell/\ell!=1$. בהזזת אינדקס מקבלים
\begin{align}
\mathbb{E}[L]&=\ee^{-\lambda}\sum_{\ell=1}^\infty
\frac{\lambda^\ell}{(\ell-1)!}=\lambda,\label{eq:he03-poisson-mean}\\
\mathbb{E}[L(L-1)]&=\ee^{-\lambda}\sum_{\ell=2}^\infty
\frac{\lambda^\ell}{(\ell-2)!}=\lambda^2,\label{eq:he03-poisson-factorial}\\
\operatorname{Var}(L)&=\mathbb{E}[L(L-1)]+\mathbb{E}[L]-\mathbb{E}[L]^2=\lambda.\label{eq:he03-poisson-var}
\end{align}
לפיכך $F=1$ בפואסון בעל תוחלת חיובית. ב-$\lambda=0$ ההתפלגות מרוכזת באפס והיחס אינו מוגדר. שוויון תוחלת ושונות הוא בדיקת עקביות שימושית למודל פואסוני, אך אינו מאפיין אותו באופן יחיד.

\section{ניסוי לוריא--דלברוק: מתי נוצרת העמידות?}\label{sec:he03-fluctuation}
מגדלים תרביות חיידקים עצמאיות ממספר התחלתי קטן, ואז חושפים אותן לבקטריופאג', נגיף שתוקף חיידקים. חיידקים רגישים נהרגים, ואילו עמידים יכולים ליצור מושבות. השאלה היא אם העמידות נוצרת בעקבות המפגש עם הנגיף, או אם מוטציות ספונטניות קדמו לחשיפה והנגיף רק בורר את העמידים שכבר היו קיימים.

במודל ההשראה הפשוט, לכל אחד מ-$N$ חיידקים סיכוי קטן וזהה $q$ לשרוד בזמן החשיפה, באופן בלתי-תלוי. מספר השורדים קרוב לפואסון עם תוחלת $Nq$ ושונות דומה. במודל המוטציות המוקדמות, לעומת זאת, אירוע מוקדם יוצר שושלת גדולה של צאצאים עמידים. לכן תרביות שונות עשויות לתת ספירות שונות מאוד, כולל זכיות נדירות וגדולות, המכונות \textenglish{jackpots}. אין להשמיט אותן רק משום שהן אינן מתאימות למודל הפואסוני.

\begin{figure}[htbp]
\centering
\begin{english}
\begin{tikzpicture}[>=stealth,every node/.style={font=\small}]
\draw[->] (0,0)--(10,0) node[right] {$n$};
\foreach \x/\lab in {0/0,2/1,4/2,6/3,8/4} {
\draw (\x,0.08)--(\x,-0.08) node[below] {$\lab$};}
\draw[->,thick,accent] (2,0.45)--(8,0.45);
\node[accent,above] at (5,0.45) {early event: $2^{4-1}=8$ descendants};
\draw[->,thick,primary] (6,1.5)--(8,1.5);
\node[primary,above] at (5,1.5) {late event: $2^{4-3}=2$ descendants};
\node[draw,anchor=west] at (8.25,0.65) {$g=4$};
\end{tikzpicture}
\end{english}
\caption{סכמה מקורית של הגברת מוטציה בזמן גידול. אירוע יחיד בדור $n$ מותיר $2^{g-n}$ צאצאים בדור המדידה $g$. זו המחשת מודל, ולא נתוני ניסוי; אירוע מוקדם נדיר יותר מפני שיש אז פחות תאים, אך השפעתו גדולה יותר.}\label{fig:he03-jackpot}
\end{figure}

\subsection{הנחות וסימון הדורות}\label{sec:he03-generations}
נניח $n_0$ תאים התחלתיים רגישים, חלוקות מסונכרנות וללא תמותה, גידול זהה של עמידים ורגישים, והורשה יציבה של המוטציה. נזניח מוטציות חוזרות באותה שושלת, מוטציות חזרה והשפעת המוטנטים על מאגר התאים הרגישים. קירוב זה מתאים לשיעור מוטציה קטן ולחלק עמידים קטן; בפרט נרצה בקירוב $g\mu\ll1$.

נשתמש \textbf{בעקביות} בדורות $n=1,\ldots,g$, כאשר הזדמנויות המוטציה נספרות בתאים בדור $n$. מספר התאים בדור זה, הגודל הסופי והפרמטר הפואסוני הם
\begin{equation}\label{eq:he03-generation-definitions}
N_n=n_0 2^n,\qquad N=n_0 2^g,\qquad
m_n\sim\operatorname{Pois}(\lambda_n),\quad \lambda_n=\mu n_0 2^n.
\end{equation}
כאן $m_n$ הוא מספר \emph{אירועי} המוטציה החדשים בדור $n$, ואילו $\mu$ הוא הסתברות קטנה לתא לכל שלב במוסכמה זו. זה אינו קצב ביחידות זמן; להמרתו לקצב בזמן נדרש זמן דור והגדרה מדויקת של הזדמנות המוטציה. נניח שה-$m_n$ בלתי-תלויים במסגרת הקירוב.

כל אירוע מותיר $a_n=2^{g-n}$ צאצאים עמידים בזמן המדידה. לכן מספר העמידים הסופי הוא
\begin{equation}\label{eq:he03-mutant-sum}
y_n=a_nm_n,\qquad a_n=2^{g-n},\qquad s=\sum_{n=1}^g a_nm_n.
\end{equation}
חשוב להבחין: $m_n$ מונה אירועים, $y_n$ מונה את צאצאיהם, ו-$s$ את כלל התאים העמידים. גם אם סכום \emph{האירועים} פואסוני, סכום הצאצאים המשוקלל אינו בדרך כלל פואסוני.

\begin{remark}
גורמים של $2$ תלויים במוסכמת הספירה: ספירת תאים לפני חלוקה, לאחריה, או ספירת חלוקות עצמן מגדירה פרמטר מוטציה אחר. כאן אין מוטנטים התחלתיים ואין איבר $n=0$. הוספת איבר כזה משנה גם את סכום האירועים וגם את המומנטים; אין לערבב גבולות סכימה שונים. הקירוב הפואסוני אף מתיר מתמטית מספר מוטציות גדול ממספר התאים, ולכן אין לפרשו כמודל מדויק מחוץ למשטר האירועים הנדירים.
\end{remark}

\section{פונקציות יוצרות: מסכומים למכפלות}\label{sec:he03-pgf}
\begin{definition}{פונקציה יוצרת הסתברות}{he03-pgf}
למשתנה שלם ואי-שלילי $X$ נגדיר
\begin{equation}\label{eq:he03-pgf}
G_X(z)=\mathbb{E}[z^X]=\sum_{k=0}^\infty P(X=k)z^k.
\end{equation}
הטור מתכנס לפחות עבור $|z|\le1$. מתקיים $G_X(1)=1$ ו-$G_X(0)=P(X=0)$, כאשר האיבר הקבוע בטור נשמר בהצבת אפס.
\end{definition}
אם המומנטים המתאימים סופיים, גזירה איבר-איבר, ובמידת הצורך גבול משמאל ב-$z=1$, נותנת
\begin{align}
G_X'(1)&=\mathbb{E}[X],\label{eq:he03-pgf-first}\\
G_X''(1)&=\mathbb{E}[X(X-1)],\label{eq:he03-pgf-second}\\
\operatorname{Var}(X)&=G_X''(1)+G_X'(1)-[G_X'(1)]^2.\label{eq:he03-pgf-var}
\end{align}
הנגזרת השנייה נותנת את $\mathbb{E}[X(X-1)]$. כדי לקבל את המומנט השני מוסיפים את התוחלת, ואז מחסרים את ריבוע התוחלת לחישוב השונות.

\begin{proposition}{סכום משתנים בלתי-תלויים}{he03-pgf-product}
אם $X,Y$ בלתי-תלויים, אז
\begin{equation}\label{eq:he03-pgf-product}
G_{X+Y}(z)=G_X(z)G_Y(z).
\end{equation}
\end{proposition}
\begin{proof}
משתמשים ב-$z^{X+Y}=z^Xz^Y$ ובפירוק תוחלת המכפלה עבור משתנים בלתי-תלויים. ללא אי-תלות אי אפשר בדרך כלל לבצע את הפירוק.
\end{proof}

למשתנה פואסוני $m_n$, סכימת טור האקספוננט נותנת
\begin{align}
G_{m_n}(u)&=\ee^{-\lambda_n}\sum_{k=0}^\infty\frac{(\lambda_nu)^k}{k!}
=\exp[\lambda_n(u-1)],\label{eq:he03-poisson-pgf}\\
G_{y_n}(z)&=G_{m_n}(z^{a_n})
=\exp[\lambda_n(z^{a_n}-1)],\label{eq:he03-weighted-pgf}\\
G_s(z)&=\prod_{n=1}^g G_{y_n}(z)
=\boxed{\exp\!\left[\sum_{n=1}^g\lambda_n(z^{a_n}-1)\right]}.\label{eq:he03-total-pgf}
\end{align}
זהו תיאור מלא של התפלגות ספירת העמידים בקירוב המודל, בלי לפרוס מפורשות את כל הסתברויותיה. בדיקת נרמול מיידית היא $G_s(1)=1$.

\begin{remark}
הגודל $G_X(z)=\mathbb{E}[z^X]$ מכונה לעיתים בדיבור בהרצאה פונקציה יוצרת מומנטים. במינוח המדויק זו פונקציה יוצרת \emph{הסתברות}. הפונקציה יוצרת המומנטים היא $M_X(t)=\mathbb{E}[\ee^{tX}]=G_X(\ee^t)$ כאשר היא קיימת. בפרט $M_X(0)=1$, בעוד $G_X(0)=P(X=0)$. ההבחנה מונעת הצבה בנקודה שגויה.
\end{remark}

\section{הסתברות לאפס עמידים ואמידת המוטציה}\label{sec:he03-mle}
אפס עמידים מתקבל אם לא קרה אף אירוע מוטציה באף דור. הצבת $z=0$ נותנת
\begin{align}
P(s=0)&=G_s(0)=\exp(-\Lambda),\qquad \Lambda=\sum_{n=1}^g\lambda_n,\label{eq:he03-zero}\\
\Lambda&=\mu n_0\sum_{n=1}^g2^n
=\mu n_0(2^{g+1}-2)=2\mu(N-n_0)\simeq2\mu N.\label{eq:he03-lambda}
\end{align}
הקירוב האחרון דורש $N\gg n_0$. $\Lambda$ הוא תוחלת מספר אירועי המוטציה, ואינו תוחלת מספר העמידים הסופי.

\begin{proposition}{אומד נראות מרבית משיעור התרביות הריקות}{he03-mle}
מתוך $c$ תרביות עצמאיות שוות בגודלן ובתנאי גידולן, נסמן ב-$c_0$ את מספר התרביות ללא עמידים. אם $0<c_0<c$, אז
\begin{equation}\label{eq:he03-mle}
\widehat\Lambda=\ln\frac c{c_0},\qquad
\widehat\mu=\frac{\ln(c/c_0)}{2(N-n_0)}\simeq\frac{\ln(c/c_0)}{2N}.
\end{equation}
\end{proposition}
\begin{proof}
לכל תרבית הסתברות $p_0=\ee^{-\Lambda}$ להיות ריקה. לכן $c_0\sim\operatorname{Bin}(c,p_0)$, והנראות, עד כדי גורם שאינו תלוי בפרמטר, היא
\begin{align}
\mathcal L(\Lambda)&\propto\ee^{-c_0\Lambda}(1-\ee^{-\Lambda})^{c-c_0},\label{eq:he03-likelihood}\\
\ell(\Lambda)&=-c_0\Lambda+(c-c_0)\ln(1-\ee^{-\Lambda}),\label{eq:he03-log-likelihood}\\
\ell'(\Lambda)&=-c_0+(c-c_0)\frac{\ee^{-\Lambda}}{1-\ee^{-\Lambda}}=0.\label{eq:he03-score}
\end{align}
כפל במכנה נותן $c_0=c\ee^{-\widehat\Lambda}$, ומכאן התוצאה. זו נקודת מקסימום כי
\begin{equation}\label{eq:he03-concavity}
\ell''(\Lambda)=-(c-c_0)\frac{\ee^{-\Lambda}}{(1-\ee^{-\Lambda})^2}<0.
\end{equation}
\end{proof}

\begin{example}{אמידה מספרית: המחשה משלימה}{he03-mutation-estimate}
אם $c=20$, $c_0=5$ ו-$N=10^9$ עם $n_0\ll N$, אז
\begin{equation}\label{eq:he03-mutation-numeric}
\widehat\Lambda=\ln4\simeq1.386,\qquad
\widehat\mu\simeq\frac{1.386}{2\cdot10^9}=6.93\times10^{-10}.
\end{equation}
זהו פרמטר לתא לשלב לפי מוסכמת המודל. האמידה משתמשת רק בשאלה אם כל תרבית ריקה, ולכן אינה נשלטת על ידי גודל הזכיות בתרביות החיוביות. מחיר הפשטות הוא ויתור על חלק מהמידע בספירות עצמן.
\end{example}

\begin{remark}
אם $c_0=c$, המקסימום הוא בגבול $\widehat\Lambda=0$; אין בכך הוכחה שקצב המוטציה אפס. אם $c_0=0$, אין אומד סופי במודל הפואסוני הזה: הנראות גדלה כאשר $\Lambda\to\infty$, והנתונים אינם נותנים אומדן עליון סופי. יש להתאים את משך הגידול או את גודל התרבית כדי למדוד גם תרביות ריקות וגם חיוביות. שיטת האפסים מניחה גם זיהוי אמין של כל תרבית שבה יש עמידים; דגימת חלק מהתרבית או יעילות גידול נמוכה משנות את הסתברות האפס.
\end{remark}

\section{השלמה מתמטית: הממוצע, השונות ומקדם פאנו}\label{sec:he03-moments}
\textbf{גבול ההרצאה:} חישוב השונות המלא נדחה להמשך. הפיתוח הבא הוא השלמה ללמידה עצמית, על סמך המודל והחוברת, ולא טענה שכל שלב הוצג בכיתה.

דרך ישירה משתמשת בליניאריות התוחלת ובאי-תלות ה-$m_n$:
\begin{align}
\mathbb{E}[s]&=\sum_{n=1}^g a_n\lambda_n
=\sum_{n=1}^g2^{g-n}\mu n_0 2^n
=g\mu N,\label{eq:he03-mutant-mean}\\
\operatorname{Var}(s)&=\sum_{n=1}^g a_n^2\lambda_n
=\mu n_0 2^{2g}\sum_{n=1}^g2^{-n}\label{eq:he03-mutant-var-start}\\
&=\mu n_0 2^{2g}(1-2^{-g})
=\mu N(2^g-1).\label{eq:he03-mutant-var}
\end{align}
כל דור תורם אותו $\mu N$ לממוצע: מיעוט האירועים המוקדמים מתאזן עם ריבוי צאצאיהם. בשונות מופיע ריבוע גודל השושלת, ולכן הדורות המוקדמים שולטים.

אפשר לקבל אותה תוצאה מהפונקציה יוצרת הקומולנטים $K_s(t)=\ln M_s(t)=\ln G_s(\ee^t)$:
\begin{align}
K_s(t)&=\sum_{n=1}^g\lambda_n(\ee^{a_nt}-1),\label{eq:he03-cgf}\\
K_s'(0)&=\sum_{n=1}^g\lambda_na_n=\mathbb{E}[s],\label{eq:he03-cgf-first}\\
K_s''(0)&=\sum_{n=1}^g\lambda_na_n^2=\operatorname{Var}(s).\label{eq:he03-cgf-second}
\end{align}
הזיהוי האחרון נובע מ-$K''=M''/M-(M'/M)^2$ וב-$t=0$ מתקיים $M=1$. ההשלמה גם מבהירה מדוע נגזרת שנייה של \emph{לוגריתם} הפונקציה יוצרת המומנטים נותנת שונות.

\begin{proposition}{חתימת הזכיות הנדירות}{he03-jackpot-fano}
עבור $g\ge1$ ו-$\mu>0$ במודל זה,
\begin{equation}\label{eq:he03-jackpot-fano}
F_s=\frac{\operatorname{Var}(s)}{\mathbb{E}[s]}
=\frac{2^g-1}{g}\simeq\frac{N}{g n_0}.
\end{equation}
במספר דורות גדול זהו יחס גדול בהרבה מ-$1$, בניגוד למודל הפואסוני הפשוט של אירועים בזמן החשיפה בלבד.
\end{proposition}
ב-$g=1$ מתקבל $s=m_1$ ולכן $F_s=1$, כנדרש. ב-$g=20$ מתקבל $F_s\simeq5.24\times10^4$: הממוצע לבדו מפספס הבדל עצום בין המודלים. היעלמות $\mu$ מהיחס אינה אומרת שהממוצע והשונות אינם תלויים בו; שניהם פרופורציונליים לו.

במספר סופי של תרביות, ייתכן שלא נדגום כלל את המוטציות המוקדמות הנדירות, ולכן הממוצע והשונות המדגמיים יכולים להשתנות מאוד בין חזרות. יתרה מכך, תנודתיות עודפת אינה הוכחה יחידה למנגנון: שונות בגודל התרביות או בקצב הגידול יכולה אף היא להגדיל את הפיזור. כוחו של הניסוי הוא בשילוב תכנון ניסויי, הורשת העמידות, והשוואה כמותית בין המודלים. גם המינוח ההיסטורי של השראה לעומת ברירה מתייחס כאן למודלים מוגדרים, ולא לשלילה של כל תגובה תאית לסביבה.

\section{סיכום ובדיקת הבנה}\label{sec:he03-recap}
משפט בייס מפריד בין הסתברות התצפית בהינתן השערה לבין הסתברות ההשערה לאחר התצפית. מתאם מתאר קשר ליניארי בלבד, ולכן יש לבדוק גם את צורת הנתונים. החלוקה הבינומית מראה כיצד רעש יכול לכייל את מספר המולקולות, והגבול הפואסוני מספק מודל ייחוס לאירועים נדירים ועצמאיים. בניסוי לוריא--דלברוק זמן האירוע קובע את גודל השושלת, ולכן מספר העמידים הוא סכום משוקלל של אירועים פואסוניים, עם תנודות גדולות במיוחד. שיעור התרביות הריקות מספק דרך לאמוד את פרמטר המוטציה בלי להסתמך על ממוצע רגיש לזכיות נדירות.

\begin{enumerate}
\item שחזרו את $23\%$ בדוגמת הבדיקה והסבירו מדוע אין להחליף רגישות בהסתברות למחלה לאחר תוצאה חיובית.
\item מדוע $\mathbb{E}[\Delta M]=0$ אינו מחייב שונות אפס? הראו היכן אנטי-המתאם בין הבנות נכנס לחישוב.
\item אם שיפוע גרף $\sigma_{\Delta Y}$ מול $\sqrt{Y_0}$ הוא $k$, מהו $\alpha$? תשובה: $k^2$.
\item הראו באמצעות פונקציות יוצרות מדוע מספר אירועי המוטציה פואסוני, אך מספר צאצאיהם בדרך כלל אינו פואסוני.
\item שחזרו את אומד שיטת האפסים, ציינו את מוסכמת הדורות, וטפלו בנפרד ב-$c_0=0$ וב-$c_0=c$.
\item כהשלמה: הסבירו מדוע כל דור תורם באותה מידה לתוחלת, אך הדורות המוקדמים תורמים יותר לשונות.
\end{enumerate}

\subsection*{קריאה נוספת}
\begin{itemize}
\item אמיר ארז, חוברת הקורס \textenglish{\emph{Physical Models of Living Systems} (2026)}, סעיף 2.8 למשפט בייס, סעיפים 2.12--2.13 למתאם והצגת נתונים, ופרק 3 להתפלגויות בדידות ולניסוי התנודות.
\item \textenglish{Philip Nelson, \emph{Physical Models of Living Systems}}: רקע למודלים הסתברותיים בביולוגיה ולפרשנות תנודות; יש להתאים את המספור למהדורה שברשות הקורא.
\end{itemize}
